\subsection{Algebraic varieties}

\bdefn

    Let $k$ be a field, then denote by $\fun(U,k)$ the $k$-algebra of set-theoretic functions $U\to k$.

    A {\emphcolor space with functions} ({\emphcolor SWF}) over $k$ is a pair $(X,\O_X)$ where $X$ is a topological space and
    for every $U\subseteq X$ open, $\O_X(U)$ is a subset of $\fun(U,k)$.
    Furthermore it must satisfy:
    \benum
        \item {\emphcolor presheaf condition}: if $f\in\O_X(U)$ and $V\subseteq U$ then the restriction of $f$ to $V$, $f|_V$, is in $\O_X(V)$.
        \item {\emphcolor sheaf condition}: if $U_i$ is an open covering of $U\subseteq X$, and $f\colon U\to k$ is a function such that
        $f|_{U_i}\in\O_X(U_i)$ for all $i$, then $f\in\O_X(U)$.
        \item {\emphcolor local condition}: for every $f\in\O_X(U)$, the subset $D_U(f)=\set{x\in U}[f(x)\neq0]$ is open, and the function
        $1/f\colon D_U(f)\to k$ is in $\O_X(D_U(f))$.
    \eenum
    $\O_X$ is called the {\emphcolor structure sheaf} of the SWF.

\edefn

\bnote

    In general a {\emphcolor presheaf} valued in a category $\C$ over a topological space $X$ is a contravariant functor
    $$ \O\colon\Op(X)^\op \to \C , $$
    where $\Op(X)$ is the category of open sets of $X$ (viewed as a preorder, ordered by inclusion).
    For $U\subseteq V$ open we denote the map $\O(V)\to\O(U)$ by $\res^V_U$, if $\C$ is a concrete category (its objects are sets) then
    we write $\res^V_U(f)=f|_U$ for $f\in\O(V)$.

    The objects $\O_X(U)$ are called {\emphcolor sections} of the presheaf.

    If $\C$ is complete (has all limits), then a presheaf $\O$ is a {\emphcolor sheaf} if for every open cover $U_i$ of $U$, the following
    is an equalizer diagram:
    $$ \O(U) \xtos{\res_0} \prod_i\O(U_i) \xtto{\res_1}[\res_2] \prod_{i,j}\O(U_i\cap U_j) , $$
    where $\res_0$ is the map $f\mapsto(f|_{U_i})$ (i.e.\ its component $\O(U)\to\O(U_i)$ is $\res^U_{U_i}$).
    The map $\res_1$ is the map $(f_i)\mapsto(f_i|_{U_i\cap U_j})$ and $\res_2$ is the map $(f_i)\mapsto(f_j|_{U_i\cap U_j})$.
    That is, their components $\prod_i\O(U_i)\to\O(U_i\cap U_j)$ are $\res^{U_i}_{U_i\cap U_j}\circ\proj_i$ and $\res^{U_j}_{U_i\cap U_j}\circ\proj_j$
    respectively.

    For a concrete category $\C$ being a sheaf means that if $f_i\in\O_X(U_i)$ are sections such that $f_i|_{U_i\cap U_j}=f_j|_{U_i\cap U_j}$ for all
    $i,j$, then there exists a unique $f\in\O_X(U)$ such that $f|_{U_i}=f_i$ for all $i$.

    Now a {\emphcolor subsheaf} of a sheaf $\O$ is another sheaf $\O'$ (on the same category) such that $\O'(U)$ is a subobject of $\O(U)$
    for all open $U$.
    Under sufficiently weak conditions (e.g.\ $\C$ being finitely complete) this is equivalent to $\O'$ being a subfunctor of $\O$
    (i.e.\ there being a monic natural transformation $\O'\To\O$).
    Conditions (1) and (2) above are equivalent to saying that $\O$ is a subsheaf of $\fun(\bul,k)$ (defined in the obvious way).

    A pair $(X,\O_X)$ where $\O_X$ is a sheaf on $X$ valued in the category of rings (resp.\ $k$-algebras) is called a {\emphcolor ringed
    space} (resp.\ a $k$-ringed space).

    Condition (3) is a replacement for the requirement that $(X,\O_X)$ be a {\emphcolor locally ringed space}.
    This condition says that the stalks of $\O_X$ (the ring of germs of a point $p\in X$, not defined here) is a local ring.

\enote

\bdefn

    A morphism of SWFs $(X,\O_X)\to(Y,\O_Y)$ is a continuous map $\phi\colon X\to Y$ such that for every open $V\subseteq Y$ and every
    $f\in\O_Y(V)$, the pullback $\phi^*(f)=f\circ\phi\colon\phi^{-1}V\to V\to k$ belongs to $\O_X(\phi^{-1}V)$.

\edefn

The identity is a morphism of SWFs because $\id^*f=f$.
And the composition of morphisms of SWFs is a morphism of SWFs clearly as well, thus SWFs form a category.

\bexam

    Let $k$ be a Hausdorff topological field, and $X$ a topological space.
    Define $(X,\O_{\it cont})$ to be the SWF where $\O_{\it cont}(U)$ is the set of continuous functions $U\to k$.
    This is easily seen to be a SWF.

    If $X$ is a real manifold, then we have SWFs $(X,\O_{C^n})$, $(X,\O_{C^\infty})$, and $(X,\O_{\it anal})$ which map an open set
    to the algebra of real valued $C^n$-, $C^\infty$-, or analytic maps respectively.
    These are SWFs and the identity maps
    $$ (X,\O_{C^n}) \xtos\id (X,\O_{C^\infty}) \xtos\id (X,\O_{\it anal}) $$
    are morphisms of SWFs.

    For a complex manifold we can define the SWF of holomorphic functions on an open subset.

\eexam

\bdefn

    If $(X,\O)$ is a space with functions and $Y\subseteq X$ is {\it open}, we can restrict $\O$ to open subsets of $Y$ to obtain another
    SWF.
    That is, define $(Y,\O|_Y)$ to be the SWF defined by $\O|_Y(U)=\O(U)$ for every open $U\subseteq Y\subseteq X$.

    In general if $Y\subseteq X$ is a subspace, we define the {\emphcolor induced space} $(Y,\O|_Y)$ as follows.
    For $V\subseteq Y$ open, let $\O|_Y(V)$ be the set of functions $f\colon V\to k$ such that there is an open cover $U_i\subseteq X$ of
    $V$ (so $V=\bigcup U_i\cap V$) and $g_i\in\O(U_i)$ such that $f|_{V\cap U_i}=g|_{V\cap U_i}$ for all $i$.

\edefn

Equivalently $f\in\O|_Y(V)$ iff for every $y\in V$ there is a neighborhood $W$ of $y$ in $X$ and a function $g\in\O_X(W)$ such that
$g|_{W\cap V}=f|_{W\cap V}$.
If we think of $\O_X$ as locally smooth maps, then this says that a map is smooth on a subspace if it can be extended locally.

\bprop

    The pair $(Y,\O|_Y)$ is an SWF and inclusion $\iota\colon X\inj Y$ is a morphism $(Y,\O|_Y)\to(X,\O_X)$.

    Furthermore if $(Z,\O_Z)$ is another SWF and $\phi\colon Z\to Y$ a set-theoretic map, then $\phi$ is a morphism of SWFs
    $(Z,\O_Z)\to(Y,\O|_Y)$ iff $\iota\circ\phi\colon(Z,\O_Z)\to(X,\O_X)$ is a morphism of SWFs.

\eprop

\bproof

    Let $V\subseteq V'\subseteq Y$ be open subsets, and $f\in\O|_Y(V')$.
    Then there is an open cover of $V'$ in $X$, $V'=\bigcup_iU_i\cap V'$ and $g_i\in\O_X(U_i)$ which restrict to $f$ on $U_i\cap V'$.
    These $U_i$ cover $V$ as well then, and $g_i$ restrict to $f|_V$ on $U_i\cap V$, so $f|_V\in\O|_Y(V)$, so it satisfies the presheaf
    condition.

    Let $V\subseteq Y$ be open and $f\in\fun(V,k)$ such that there is an open cover $V=\bigcup_iV_i$ with $f|_{V_i}\in\O_Y(V_i)$.
    Then we cover each $V_i$ by $U^j_i$ such that there are $g^j_i\in\O_X(U^j_i)$ with $g^j_i|_{U^j_i\cap V_i}=f|_{V_i\cap U^j_i}$.
    But then the collection of all $U^j_i$ as $i,j$ ranges and $g^j_i$ show that $f\in\O|_Y(V)$.
    So $\O|_Y$ is a sheaf.

    Now we need to show the locality axiom.
    Let $V\subseteq Y$ be open and $f\in\O|_Y(V)$ so there are $U_i\subseteq X$ open covering $V$ such that $f|_{U_i\cap V}=g_i|_{U_i\cap V}$
    for $g_i\in\O_X(U_i)$.
    Then
    $$ D_V(f) = \set{x\in V}[f(x)\neq0] = \bigcup_i\set{x\in V\cap U_i}[f(x)\neq0] = \bigcup_i D_{U_i}(g_i)\cap V , $$
    which is relatively open in $V$.
    And we see the $D_{U_i}(g_i)$ cover $D_V(f)$ and $1/f$ restricts to $1/g_i\in\O_X(D_{U_i}(g_i))$ their common domain, so $1/f\in\O|_Y(D_U(f))$.
    Thus $\O|_Y$ is indeed a SWF.

    The inclusion is continuous by virtue of the subspace topology.
    Let $U\subseteq X$ open so $\iota^{-1}U=U\cap Y$, and let $f\in\O_X(U)$ then $\iota^*(f)=f\circ\iota=f|_{U\cap V}$.
    This is in $\O|_Y(U\cap Y)$ by considering the open cover $U$ and $f\in\O_X(U)$.

    The composition of morphisms of SWFs is a morphism of SWFs, so if $\phi$ is a morphism so is $\iota\circ\phi$.
    Conversely if $\iota\circ\phi$ is continuous, then so is $\phi$ by virtue of the subspace topology.
    And if $f\in\O|_Y(V)$ for $V\subseteq Y$ open, then let $g_i\in\O_X(U_i)$ such that $g_i|_{U_i\cap V}=f|_{U_i\cap V}$ with $U_i$ covering
    $V$.
    Then
    $$ (\iota\circ\phi)^*(g_i)=g_i\circ\iota\circ\phi=g_i|_{U_i\cap V}\circ\phi=f|_{U_i\cap V}\circ\phi=f\circ\phi|_{W_i} $$
    is in $\O_X(W_i)$ with $W_i=(\iota\circ\phi)^{-1}(U_i)$.
    These $W_i$ cover $\phi^{-1}V$ and thus by definition $f\circ\phi\in\O|_Y(\phi^{-1}V)$, showing $\phi$ is a morphism of SWFs.
    \qed

\eproof

\bexam

    Let $B^n\subseteq\bR^n$ or $\bC^n$ be the open unit ball.
    Then the above proposition gives us SWFs
    $$ (B^n,\O_{\it cont}),(B^n,\O_{C^k}),(B^n,\O_{C^\infty}),(B^n,\O_{\it anal}),(B^n,\O_{\it holo}) . $$
    We can then define a real/complex topological/$C^k$/$C^\infty$/analytic manifold as a SWF $(X,\O)$ such that $X$ is Hausdorff and second
    countable and there is an open cover $X=\bigcup_iU_i$ such that there is a single SWF from one of the above ($(B_n,\O_\bul)$) such that
    each $i$ $(U_i,\O|_{U_i})$ is isomorphic to it.

    This is an equivalent definition to the standard definition of a manifold via charts.
    Then a smooth map between manifolds is just a morphism of SWFs.

\eexam

We use this example to define general algebraic varieties.
We will define affine algebraic varieties as a special family of SWFs, and then a general algebraic variety will be a SWF locally isomorphic
to an affine algebraic variety.

\bdefn

    Let $Y\subseteq\bA^n$ be a subspace.
    A function $f\colon Y\to k$ is {\emphcolor regular} when it is locally rational (the product of polynomials).
    That is, for every $P\in Y$ there is a neighborhood $U_P\subseteq Y$ of $Y$ and polynomials $g_P,h_P\in k[\bA^n]$ such that
    $h_P(Q)\neq0$ for $Q\in U_P$ and $f(Q)=g_P(Q)/h_P(Q)$ for all $Q\in U_P$.

    Denote the set of regular functions $Y\to k$ by $\O_\reg(Y)$.

\edefn

\blemm

    The pair $(Y,\O_\reg)$ is a SWF for every $Y\subseteq\bA^n$.

\elemm

\bproof

    It is easy to see that $\O_\reg(U)\subseteq\fun(U,k)$ is a subalgebra; i.e.\ that it is closed under addition and multiplication,
    and contains the unit (and thus all constants).
    Furthermore since being regular is a local condition, it is again easy to see that $\O_\reg$ is a sheaf.

    We just need to show locality.
    Let $f\in\O_\reg(U)$ then for $P\in U$ there is an open $P\in U_P\subseteq U$ such that $f=g_P/h_P$ for polynomials $g_P,h_P$
    in $U_P$.
    Then
    $$ D_U(f) = \set{x\in U}[f(x)\neq0] = \bigcup_{P\in U}\set{x\in U_P}[g_P(x)\neq0] = \bigcup_{P\in U_P}D_{U_P}(g_P) $$
    is open since each $g_P$ is continuous.
    And on $D_{U_P}(g_P)$ we have $1/f=h_P/g_P$, so $1/f$ is regular on $D_U(f)$ as desired.
    \qed

\eproof

\bnote

    It is easy to see that $\O_\reg$ is the smallest SWF containing all the polynomial functions in $k[\bA^n]$.
    Such a SWF must contain regular functions because of locality.

\enote

\blemm

    Let $U\subseteq\bA^n$ be open, and $f\in\O_\reg(U)$ a regular function.
    Then $f$ is rational: there are $g,h\in k[\bA^n]$ such that $h\neq0$ on $U$ and $f=g/h$.

\elemm

\bproof

    By definition there is an open covering $U=\bigcup_iU_i$ and polynomials $g_i,h_i$ such that $h_i\neq0$ on $U_i$
    and $f|_{U_i}=g_i/h_i$ on $U_i$.
    Further we can assume that $g_i,h_i$ are relatively prime and $U_i$ are nonempty.

    We claim that for any $i$, $h_i$ is nonzero on all of $U$, and $f=g_i/h_i$ on all of $U$.
    Let $P\in U_j$, so we have
    $$ \lbar{\frac{g_i}{h_i}}_{U_i\cap U_j} = f|_{U_i\cap U_j} = \lbar{\frac{g_j}{h_j}}_{U_i\cap U_j} . $$
    Thus $g_ih_j-g_jh_i=0$ on $U_i\cap U_j$.
    Since $\bA^n$ is irreducible, any nonempty open subset is dense (this is equivalent) and so the intersection of nonempty open subsets
    is then dense (as each is dense and thus their intersection is nonempty, and thus a nonempty open subset, so dense).
    Thus $g_ih_j$ and $g_jh_i$ agree on a dense set, so they are equal (two maps into a Hausdorff space agreeing on a dense set are equal).

    Since $k[\bA^n]$ is a UFD and $\set{g_i,h_i}$ are all relatively prime, we must have $h_i=ch_j$ and $g_i=cg_j$ for some
    $c\in k[\bA^n]^\times$ (which is $k^\times$).
    Thus we have $h_i(P)=c(P)h_j(P)\neq0$ for all $P\in U_j$, so $h_i$ is nonzero on all of $U$.
    And so for $P\in U_j$,
    $$ \frac{g_i(P)}{h_i(P)} = \frac{g_j(P)}{h_j(P)} = f(P) , $$
    so that $f=g_i/h_i$ on all of $U$ as desired.
    \qed

\eproof

So $(\bA^n,\O_\reg)$ is a much simpler object than the general case.

\bdefn

    An {\emphcolor affine algebraic variety} (or {\emphcolor affine variety}) is a SWF $(X,\O_X)$ which is isomorphic to a SWF
    $(Y,\O_\reg)$ where $Y\subseteq\bA^n$ is closed.

    An {\emphcolor algebraic variety} is a SWF $(X,\O_X)$ which has a finite local covering $X=\bigcup_{i=1}^nU_i$ such that each
    $(U_i,\O_X|_{U_i})$ is an affine variety.
    Elements of a section of an algebraic variety (elements of $\O_X(U)$) are called {\emphcolor regular functions}.

    A morphism of algebraic varieties is just a morphism of SWFs.

\edefn

\blemm

    Every algebraic variety is Noetherian (as a topological space).

\elemm

\bproof

    First we show that $\bA^n$ is Noethrian.
    Indeed, if $\set{Z_i}_i$ is a collection of closed subsets of $\bA^n$ then $\set{\I(Z_i)}_i\subseteq k[\bA^n]$ has a maximal element
    since $k[\bA^n]$ is Noetherian.
    Since $Z\circ\I$ is the identity on closed sets, this means there is a maximal $Z_i$, and thus $\bA^n$ is Noetherian.

    An affine variety is homeomorphic to a closed subset of $\bA^n$ which is Noetherian.

    Thus an algebraic variety is locally Noetherian, i.e.\ there is a finite open covering $\bigcup_{i=1}^nU_i$ where each $U_i$ is Noetherian.
    Then $X$ is Noetherian.
    \qed

\eproof

\bcoro

    Every algebraic variety is quasicompact.

\ecoro

\bdefn

    Let $Y\subseteq\bA^n$ be closed, then define $k[Y]=k[\bA^n]/\I(Y)$.
    Then there is a natural homomorphism
    $$ \ev_Y \colon k[\bA^n] \xtos\ev \fun(\bA^n,k) \xtos\res \fun(Y,k) . $$
    By definition the kernel of this is $\I(Y)$, and so $\ev_Y$ induces an injection $k[Y]\inj\fun(Y,k)$.

\edefn

This notation is compatible with our previous definition of $k[\bA^n]$.
This is because $\I(\bA^n)=0$ so $k[\bA^n]/\I(\bA^n)=k[\bA^n]$.

Note that $k[Y]$ is a fg $k$-algebra.
Furthermore it is reduced, as it can be embedded into $\fun(Y,k)$ which is reduced as $k$ is (also because $\I(Y)$ is a radical ideal).
Furthermore, $k[Y]$ is a domain iff $\I(Y)$ is prime iff $Y$ is irreducible.

If $B$ is a reduced fg $k$-algebra, then $B\cong k[Y]$ for some closed $Y\subseteq\bA^n$.
Indeed, $B\cong k[\bA^n]/I$ for some ideal $I$, and since $B$ is reduced $I$ is a radical ideal and thus equal to $\I(Y)$ for $Y=Z(\I)$,
so $B\cong k[Y]$.

For $f\in k[\bA^n]$ let $\bar f=f+\I(Y)\in k[Y]$ denote its projection.
Then viewed as functions $\bA^n\to k$ and $Y\to k$ respectively, we have the equality
$$ D_Y(\bar f) = Y\cap D(f) . $$
Since the $D(f)$ form a base of open subsets of $k[\bA^n]$, we then see that $\set{D_Y(\bar f)}_{\bar f\in k[Y]}$ is a basis of $Y$.

\blemm

    For every collection $\set{\bar f_i}_i\subseteq k[Y]$, $\set{D_Y(\bar f_i)}_i$ form an open cover of $Y$ iff the ideal $(\set{f_i}_i)$
    is $k[Y]$.

\elemm

\bproof

    We have that $D_Y(\bar f_i)$ cover $Y$ iff $D(f_i)$ cover $Y$.
    Since $D(f_i)=Z(f_i)^c$, this is iff $\bigcap_iZ(f_i)\cap Y=\emptyset$, equivalently $Z((f_i)_i)\cap Y$ is empty.
    But
    $$ Z((f_i)_i) \cap Y = Z((f_i)_i) \cap Z(\I(Y)) = Z((f_i)_i + \I(Y)) , $$
    and this is empty iff $Z((f_i)_i)+\I(Y)$ is all of $k[\bA^n]$ (by Nullstellensatz).
    Quotienting by $\I(Y)$ this is equivalent to $Z((\bar f_i)_i)=k[Y]$.
    \qed

\eproof

Recall that $f\in\O_\reg(Y)$ iff there is an open cover $Y=\bigcup_iU_i$ such that $f|_{U_i}$ is rational for each $i$.
Since $Y$ is quasicompact, we can assume this covering is finite, and taking basic open sets, that $U_i=D_Y(t_i)$ for $t_i\in k[Y]$.
So $f=g_i/h_i$ for $g_i,h_i\in k[\bA^n]$ and $h_i$ nonzero.
Equivalently, $f=\bar g_i/\bar h_i$ (as $\bar g$ agrees with $g$ on $Y$).

\bthrm[title=Serre]

    For every closed $Y\subseteq\bA^n$ there is an equality $k[Y]=\O_\reg(Y)$ (viewing $k[Y]$ as a subalgebra of $\fun(Y,k)$).

\ethrm

\bproof

    Since $\O_\reg(Y)$ consists of all polynomial functions, clearly $k[Y]\subseteq\O_\reg(Y)$.
    By the above if $f\in\O_\reg(Y)$ then there is a finite set of triples $\set{t_i,g_i,h_i}$ of polynomials in $k[Y]$ such that
    $D_Y(t_i)$ cover $Y$ and $f_i=g_i/h_i$ on $D_Y(t_i)$ and $h_i\neq0$ on $D_Y(t_i)$.

    By assumption we have that $D_Y(t_i)\subseteq D_Y(h_i)$, and so $D_Y(t_i)=D_Y(t_i)\cap D_Y(h_i)=D_Y(t_ih_i)$.
    Since $g_i/h_i=t_ig_i/t_ih_i$, we can replace our collection with $\set{t_ih_i,t_ig_i,t_ih_i}_i$.
    In particular we can assume that $t_i=h_i$.

    Now we want to show that $f\in k[Y]$.
    We first give a short, but tricky, proof and afterward we will give a more insightful one.
    Notice that for every $i$ we have
    $$ fh_i^2 = g_ih_i \in \fun(Y,k), $$
    so $(fh_i-g_i)h_i=0$.
    This is true in $D_Y(h_i)$ clearly, and for $P\in Y-D_Y(h_i)$ we have that $h_i(P)=0$ so the equality holds trivially.

    Now since $Y=\bigcup_iD_Y(h_i)=\bigcup_iD_Y(h_i^2)$ we conclude that $(h_i^2)=k[Y]$ as per the previous lemma.
    So there are $r_i\in k[Y]$ such that
    $$ \sum_i r_ih_i^2 = 1 . $$
    (The algebraic analogue of a partition of unity.)
    Thus
    $$ f = f\sum_i r_ih_i^2 = \sum_i r_ifh_i^2 = \sum_i r_ig_ih_i \in k[Y] $$
    as desired.
    \qed

\eproof

We now give a more insightful proof of Serre's theorem.

\blemm

    Let $A$ be a ring and suppose $f_1,\dots,f_n$ generate the unit ideal.
    Then
    $$ 0 \to A \xtos\phi \prod_i A_{f_i} \xtos\psi \prod_{i,j}A_{f_if_j} , $$
    is exact, where $\phi(a)=(a/1)_i$ and $\psi((a_i)_i)=(a_i-a_j)_{ij}$.

\elemm

\bproof

    First we show that the sequence is exact at $A$, meaning $\phi$ is injective.
    If $a/1=0$ for all $i$, then there are $n_i$ such that $f_i^{n_i}a=0$ for each $i$.
    But
    $$ f_i \in \sqrt{(f_i^{n_i})} , $$
    and so $\sqrt{(f_i^{n_i})}$ is equal to the unit ideal.
    But this means that $(f_i^{n_i})$ is equal to the unit ideal, and thus $a=0$.

    Now we show that $\im\phi=\ker\psi$.
    Firstly,
    $$ \psi\circ\phi(a) = \psi((a/1)_i) = (a/1-a/1)_{i,j} = 0 , $$
    so the sequence is a chain complex.
    We now need only show that $\ker\psi\subseteq\im\phi$.
    So let $(a_i/f_i^{n_i})_i$ lie in $\ker\psi$.
    Since $a_if^{r_i}/f_i^{n_i+r_i}=a_i/f^{n_i}_i$, we may enlarge the $n_i$ so that they are equal.

    So
    $$ 0 = \psi(a_i/f_i^n) = (a_i/f_i^n - a_j/f_j^n)_{i,j} . $$
    This means that $a_i/f_i^n=a_j/f_j^n$ in $A_{f_if_j}$ for all $i,j$ and thus there is an $M_{ij}$ such that
    $(f_if_j)^{M_{ij}}(f_j^na_i-f_i^na_j)=0$.
    Again by potentially enlarging the $M_{ij}$ we can assume they are all equal to $M$.
    So we have that for all $i,j$,
    $$ (f_if_j)^M(f_j^na_i - f_i^na_j) = 0 . $$
    As before $f_j^{M+n}$ generate the unit ideal so there are $b_j\in A$ such that
    $$ \sum_j b_jf_j^{M+n} = 1 , $$
    and so if we define $a=\sum_jb_jf_j^Ma_j$, then
    $$ f_i^{M+r}a = \sum_j b_jf_j^Mf_i^{M+r}a_j = \sum_j b_jf_i^Mf_j^{M+r}a_i = f_i^Ma_i\sum_jb_jf_j^{M+r} = f_i^Ma_i . $$
    Thus $a/1=a_i/f_i^n$, and so $\phi(a)=(a_i/f_i^n)_i$ as desired.
    \qed

\eproof

Now in the proof of Serre's theorem, since $(h_i)_i=k[Y]$ we have the exact sequence
$$ 0 \to k[Y] \xtos\phi \prod_ik[Y]_{h_i} \xtos\psi \prod_{i,j}k[Y]_{h_ih_j} . $$
Now consider $\bar f=(g_i/h_i)\in\prod_ik[Y]_{h_i}$, then $g_i/h_i-g_j/h_j=0$ in $k[Y]_{h_ih_j}$ because as we will show,
$$ (g_ih_j - g_jh_i)h_ih_j = 0 \hbox{ in $k[Y]$}. $$
Since $k[Y]\subseteq\fun(Y,k)$ we can show that the above expression vanishes for every $P\in Y$, and this is sufficient.
If $P\in D_Y(h_i)\cap D_Y(h_j)$ then
$$ (g_i/h_i)(P) = (g_j/h_j)(P) = f(P) , $$
and the equality follows.
For all other $P$ either $h_i(P)$ or $h_j(P)$ is zero, and the equality follows trivially.

Thus by exactness there is a $\tilde f\in k[Y]$ such that $\phi(\tilde f)=\bar f$, meaning $\tilde f/1=g_i/h_i\in k[Y]_{h_i}$ for all $i$.
So for some $n$, $h_i^n(h_i\tilde f-g_i)=0$
In particular for $P\in D_Y(h_i)$ since $h_i(P)\neq0$ we have that $h_i(P)\tilde f(P)=g_i(P)$ so
$$ \tilde f(P) = g_i(P)/h_i(P) = f(P) . $$
Thus $f=\tilde f\in k[Y]$ as desired, showing that $\O_\reg(Y)=k[Y]$ as desired.

\bnote

    Abusing terminology, for {\it any\/} SWF $(X,\O_X)$ we call elements of $\O_X(U)$ {\emphcolor regular functions} on $U$.

\enote


\bthrm

    Let $(X,\O_X)$ be a SWF and $Y\subseteq\bA^n$ a subset.
    Let $\psi\colon X\to Y$ be a map of sets, and set $\psi_i=\proj_i\circ\psi\in\fun(X,k)$ to be the projections of $\psi$.
    Then $\psi$ is a morphism of SWFs $(X,\O_X)\to(Y,\O_\reg)$ iff $\psi_i\in\O_X(X)$ for all $i$.
    Equivalently, $\psi^*\colon\fun(Y,k)\to\fun(X,k)$ maps $\O_Y(Y)$ into $\O_X(X)$.

\ethrm

\bnote

    Equivalently, $\psi=(\psi_1,\dots,\psi_n)\colon X\to Y$ is a morphism iff each $\psi_i$ is regular.

\enote

\bproof

    Notice that $\proj_i\colon\bA^n\to k$ is given by a regular function $x_i$ on $\bA^n$.
    So if $\psi$ is a morphism, then every $\psi_i=\psi^*(x_i)$ is a regular function on $X$.

    Conversely assume $\psi_i\in\O_X(X)$ for all $i$.
    So we want to show that $\psi$ is continuous, and for every $f\in\O_Y(U)$ we have $\psi^*(f)\in\O_X(\psi^{-1}U)$ for all $U\subseteq Y$
    open.

    First we show that for every $f\in k[\bA^n]$ we have $\psi^*(f)\in\O_X(X)$.
    Notice that $\psi^*\colon\fun(\bA^n,k)\to\fun(X,k)$ is a $k$-algebra morphism, while $f$ belongds to the $k$-subalgebra
    generated by the $x_i$.
    Thus $\psi^*f$ belongds to the $k$-algebra generated by the $\psi^*(x_i)=\psi_i$.
    Since $\psi_i\in\O_X(X)$ by assumption, this means $\psi^*f\in\O_X(X)$ as desired (since it is a subalgebra).

    Now we show that $\psi$ is continuous.
    The basic open sets of $Y$ are $D_Y(f)$ with $f\in k[\bA^n]$, and so it is sufficient to show $\psi^{-1}D_Y(f)$ is open in $X$.
    Notice that
    $$ \psi^{-1}D_Y(f) = \set{x\in X}[f(\psi x)=0] = D_X(\psi^*f) . $$
    And since $\psi^*f\in\O_X(X)$ by locality of SWFs we have that $D_X(\psi^*f)$ is open as desired.

    Now we need to show that for every open $U\subseteq Y$ and $f\in\O_Y(U)$ we have $\psi^*f\in\O_X(\psi^{-1}U)$.
    Replacing $Y$ by $U$ and $X$ by $\psi^{-1}U$, we can assume $U=Y$.

    First assume that $f$ is rational, so $f=g/h$ on $Y$ with $g,h\in k[\bA^n]$ and $h$ nonzero.
    Then $\psi^*f=\psi^*g/\psi^*h$ and $\psi^*h(x)\neq0$ for every $x\in X$.
    By locality $1/\psi^*h$ is in $\O_X(X)$ and thus $\psi^*f$ is in $\O_X(X)$ as it is a $k$-algebra, as desired.

    Now in general, locally $f\in\O_Y(Y)$ has the form $g/h|_Y$, the general case follows from the rational case and
    the sheaf axioms.

    And the final equivalence is because if $\psi^*$ maps $\O_Y(Y)$ into $\O_X(X)$ then it maps $x_i\in\O_Y(Y)$ to $\psi^*(x_i)=\psi_i$
    into $\O_X(X)$.
    \qed

\eproof

\bcoro

    For every SWF $X=(X,\O_X)$ there is a natural bijection
    $$ \hom_\SWF(X,\bA^n) \xtos\sim \O_X(X)^n,\qquad \psi\mapsto(\proj_1\circ\psi,\dots,\proj_n\circ\psi) . $$
    In particular there is a natural identification $\hom_\SWF(X,\bA^1)\cong\O_X(X)$.

    Furthermore, for $X=(X,\O_X)$ and a subset $Y\subseteq\bA^n$, a set-theoretic map $\psi\colon X\to Y$ is a morphism of SWFs
    $(X,\O_X)\to(Y,\O_\reg)$ iff the composition $X\xtos\psi Y\inj\bA^n$ is a morphism of SWFs.

\ecoro

\bproof

    The first point folows from the above theorem with $Y=\bA^n$.
    The second follows from the theorem as well, since $\psi=(\psi_1,\dots,\psi_n)$ being a morphism is equivalent to all
    $\psi_i$ being regular which is equivalent to the composition of $\psi$ with the inclusion being a morphism.
    \qed

\eproof

\bnote

    The second point follows directly from the fact that $(Y,\O_\reg)$ is the restriction of $(\bA^n,\O_\reg)$ to $Y$.

\enote

Because for an affine scheme $(X,\O_X)$ by Serre we have an identifcation $\O_X(X)\cong k[X]$, the following is immediate.

\bcoro

    If $(X,\O_X)$ and $(Y,\O_Y)$ are affine varities, then $\psi\colon X\to Y$ is a morphism of SWFs iff it maps $k[Y]$ to $k[X]$.

\ecoro

There is an obvious functor $\SWF_k^\op\to\Alg_k$ given by mapping $(X,\O_X)$ to $\O_X(X)$ and $\phi\colon X\to Y$ to $\phi^*\colon\O_Y(Y)\to\O_X(X)$.
This is called the {\emphcolor global sections} functor, as it takes global sections.
It is denoted by $\Gamma$, so $\Gamma(X,\O_X)$ for objects (we may omit one of the two inputs when the other is obvious), and $\Gamma(\phi)=\phi^*$ on morphisms.

\blemm

    A morphism of SWFs $\psi\colon X\to\bA^n$ satisfies $\phi(X)\subseteq Y$ for closed $Y\subseteq\bA^n$ iff $\psi^*\colon k[\bA^n]\to\O_X(X)$ satisfies $\I(Y)\subseteq\ker\psi^*$.

\elemm

\bproof

    We have $\psi(X)\subseteq Y$ iff $\psi(x)\in Z(\I(Y))$ for all $x\in X$ and $f\in\I(Y)$, which is iff $f(\psi(x))=\psi^*(f)(x)=0$.
    So $\psi(X)\subseteq Y$ iff $\psi^*(f)=0$ for all $f\in\I(Y)$, i.e.\ iff $\I(Y)\subseteq\ker\psi^*$.
    \qed

\eproof

\bprop

    Let $(X,\O_X)$ be an SWF and $(Y,\O_Y)$ be an affine variety.
    Then $\Gamma$ is bijective
    $$ \Gamma\colon\hom_\SWF(X,Y)\to\hom_\Alg(\O_Y(Y),\O_X(X)) . $$

\eprop

\bproof

    First take $Y=\bA^n$, so $\O_Y(Y)=k[\bA^n]$, the free $k$-algebra on $n$ generators.
    Thus there is a natural bijection
    $$ \nat\colon\hom_\Alg(k[\bA^n],\O_X(X)) \to \O_X(X)^n,\qquad \phi\mapsto(\phi(x_1),\dots,\phi(x_n)) . $$
    Now, the composition
    $$ \hom_\SWF(X,\bA^n) \xtos\Gamma \hom_\Alg(k[\bA^n],\O_X(X)) \xtos\nat \O_X(X)^n $$
    maps $\phi$ to $(\phi^*(x_1),\dots,\phi^*(x_n))$, and as we showed above this is a bijection.
    Thus $\Gamma$ must be a bijection, as desired.

    Now let $Y\subseteq\bA^n$ be an arbitrary affine variety, so $\O_Y(Y)=k[Y]=k[\bA^n]/\I(Y)$.
    Recall that $(Y,\O_Y)$ is induced from $(\bA^n,\O_\reg)$ so SWF morphisms into $Y$ are the same as SWF morphisms into $\bA^n$ landing in $Y$.
    Thus by the previous lemma
    $$ \hom_\SWF(X,Y) \cong \set{\phi\in\hom_\SWF(X,\bA^n)}[\phi(X)\subseteq Y] = \set{\phi\in\hom_\SWF(X,\bA^n)}[\I(Y)\subseteq\ker\psi^*] , $$
    and as we just showed, this is bijective to
    $$ \cong \set{\phi^*\in\hom_\Alg(k[\bA^n],\O_X(X))}[\I(Y)\subseteq\ker\psi^*] , $$
    and these are precisely the maps which quotient to $k[Y]=k[\bA^n]/\I(Y)\to\O_X(X)$, i.e.\ this forms a bijection
    $$ \cong \hom_\Alg(k[Y],\O_X(X)) . $$
    Now it is easy to see that the composition of these isomorphisms is $\Gamma$.
    \qed

\eproof

\bcoro

    A SWF $Y$ is an affine variety iff the $k$-algebra $\O_Y(Y)$ is fg and for every SWF $X$, $\Gamma$ is bijective from $\hom_\SWF(X,Y)$.

\ecoro

\bproof

    We just showed for an affine variety, the property holds.
    Conversely, if $\O_Y(Y)$ is fg then it is equal to $\O_{Y'}(Y')$ for some affine variety $Y'$ (we showed that reduced fg $k$-algebras $B$ are all of the form $k[Y']$ for an
    affine $Y'$, and $\O_Y(Y)\subseteq\fun(Y,k)$ is reduced always).
    Thus there is an isomorphism
    $$ \hom_\SWF(X,Y') \cong \hom_\Alg(\O_{Y'}(Y'),\O_X(X)) = \hom_\Alg(\O_Y(Y),\O_X(X)) \cong \hom_\SWF(X,Y) , $$
    where the final isomorphism is given by $\Gamma$.
    Thus $\hom(\bul,Y')\cong\hom(\bul,Y)$ naturally, so by Yoneda $Y'\cong Y$.
    \qed

\eproof

\bcoro

    $\Gamma$ forms an anti-equivalence of categories from affine varieties to fg reduced $k$-algebras: $\AVar_k\to\Alg_k^{\fg,\red}$.

    In particular $\Gamma$ reflects isomorphisms, so $X\cong Y$ iff $\O_X(X)\cong\O_Y(Y)$ for affine $X,Y$.

\ecoro

\bproof

    We showed that $\Gamma$ is fully faithful on affine varieties, and so it is an equivalence to its essential image.
    Its essential image are the fg reduced $k$-algebras.

    We can explicitly construct its inverse: for a fg reduced $k$-algebra $A$, then we map it to $\specm(A)$, and maps $\phi\colon A\to B$ are mapped to $\abs{\phi^*}$.
    \qed

\eproof

\bcoro

    Let $X\subseteq\bA^n$ and $Y\subseteq\bA^m$ be affine varieties.
    A map $\phi\colon X\to Y$ is a morphism of varieties iff $\phi=(f_1,\dots,f_m)$ for some $f_i\in k[\bA^n]$.

\ecoro

\bproof

    $\phi$ is a morphism iff it is a morphism into $\bA^m$, which is iff $\phi$'s projections are in $\O_X(X)=k[X]$, which are just the restriction of polynomial functions to $X$.
    \qed

\eproof

\bdefn

    Let $X$ be a topological space.
    A subset $Y\subseteq X$ is {\emphcolor locally closed} iff $Y=U\cap F$ for open $U\subseteq X$ and closed $F\subseteq X$.

\edefn

The following lemma better explains the reasoning for the terminology of ``locally closed''.

\blemm

    For $Y\subseteq X$ a subspace, the following are equivalent:
    \benum
        \item $Y$ is locally closed,
        \item $Y$ is an open subset of a closed subset,
        \item $Y$ is a closed subset of an open subset,
        \item for every $y\in Y$, there is a neighborhood $U$ of $y$ in $X$ such that $Y\cap U$ is closed in $U$.
    \eenum

\elemm

\bproof

    Recall that open (resp.\ closed) subsets of a subspace $S$ are of the form $S\cap U$ for $U$ open (resp.\ $S\cap F$ for $F$ closed).
    So an open subset of a closed subset $F$ are the sets of the form $U\cap F$ for $U$ open, and similarly for closed subsets of open subsets.
    Thus we see the equivalence of (1), (2), and (3).

    Now if $Y$ is locally closed, then it is the closed subset of an open subset $U$.
    Thus every $y\in Y$ has neighborhood $U$ in which $Y$ is closed, showing (4).
    Conversely, for each $y\in Y$ let $y\in U_y$ be a neighborhood where $Y\cap U_y$ is closed in $U_y$.
    Let $U=\bigcup_yU_y$, so
    $$ U - Y = \bigcup_y(U_y-Y) , $$
    and since $Y\cap U_y$ is closed in $U_y$, we have $U_y-Y=U_y-Y\cap U_y$ is open, so $U-Y$ is open and thus $Y$ is closed in $U$.
    \qed

\eproof

\bcoro

    A locally closed subset of a locally closed subset is locally closed.

\ecoro

\bproof

    So suppose $Z\subseteq Y\subseteq X$ is a chain of inclusions, where each subset is locally closed in the next.
    It is sufficient to show that an open/closed subset of $Y$ is locally closed in $X$, since if $Z=U\cap F$ in $Y$, then $U,F$ are locally closed in $X$ and then clearly
    their intersection is too.
    But clearly by (2) an open subset of $Y$ is locally closed, and by (3) the same for a closed subset.
    \qed

\eproof

\bprop

    Let $X\subseteq\bA^n$ be an affine variety, and let $f\in k[\bA^n]$.
    Then $(D_X(f),\O_\reg)$ is an affine variety isomorphic to $(Y,\O_\reg)$ where $Y\subseteq\bA^{n+1}$ is the closed subset $Z(\I(X)\cup\set{fx_{n+1}-1})$.

\eprop

\bproof

    We construct explicit inverse morphisms $\phi\colon D_X(f)\to Y\colon\psi$.
    Define $\phi$ by mapping $(x_1,\dots,x_n)\in X\subseteq\bA^n$ to
    $$ \phi(x_1,\dots,x_n) = \parens{x_1,\dots,x_n,\frac1{f(x_1,\dots,x_n)}} . $$
    Its image lies in $Y$ clearly, and its components are rational functions on $D_X(f)$ so it forms a morphism into $\bA^n$, thus into $Y$.
    Similarly define $\psi$ by mapping $(x_1,\dots,x_{n+1})\in Y$ to
    $$ \psi(x_1,\dots,x_{n+1}) = (x_1,\dots,x_n) . $$
    Its image lies in $X$ since $f(x_1,\dots,x_n)x_{n+1}=1$ so $f(x_1,\dots,x_n)\neq0$, and it is a morphism into $\bA^n$ and thus into $X$.
    Clearly $\psi\circ\phi=\id$ and for $(x_1,\dots,x_{n+1})\in Y$ we must have $x_{n+1}=1/f(x_1,\dots,x_n)$ by definition so $\phi\circ\psi=\id$.
    \qed

\eproof

\bexam

    The SWF $(D_{\bA^1}(x),\O_\reg)$ is isomorphic to the hyperbola $Z(xy-1)\subseteq\bA^2$.

\eexam

We can also prove the above proposition via Yoneda.
Define $B=k[X]_f$.
Now we know that $\Gamma$ defines a bijection $\hom_\SWF(Z,X)\to\hom_\Alg(k[X],\O_Z(Z))$ for any SWF $Z$.
Furthermore, the universal property of localization says that a morphism $\zeta\colon k[x]\to\O_Z(Z)$ factors through $B$ iff $\zeta(f)\in\O_Z(Z)^\times$, and such a factorization is unique.
And by locality, $\O_Z(Z)^\times$ are just the nowhere vanishing functions in $\O_Z(Z)$.
Thus $\Gamma$ restricts to a bijection
$$ \hom_\Alg(k[X]_f,\O_Z(Z)) \cong \set{\phi\in\hom_\SWF(Z,X)}[\phi^*(f)(z)\neq0\hbox{ for all $z\in Z$}] . $$
And we know
$$ \phi^*(f)(z) = f(\phi(z)) , $$
so $\phi^*(f)(z)\neq0$ iff $\phi(z)\in D_X(f)$, thus
$$ \hom_\Alg(k[X]_f,\O_Z(Z)) \cong \set{\phi\in\hom_\SWF(Z,X)}[\phi(Z)\subseteq D_X(f)] . $$
And since $D_X(f)\subseteq X$ has the induced SWF structure (as they are both induced from $\bA^n$), this means $\Gamma$ defines a natural isomorphism
$$ \hom_\Alg(k[X]_f,\O_Z(Z)) \cong \hom_\SWF(Z,D_X(f)) . $$
Now $k[X]_f$ is a reduced fg $k$-algebra, and thus it is equivalent to some affine variety $Y$ under $\Gamma$ (so $\O_Y(Y)=k[X]_f$), so we have natural isomorphisms
$$ \hom_\SWF(Z,Y) \cong \hom_\Alg(k[X]_f,\O_Z(Z)) \cong \hom_\SWF(Z,D_X(f)) , $$
and by Yoneda $Y\cong D_X(f)$ as desired.

Notice that this means $k[D_X(f)]\cong k[X]_f$ (the left-hand side is $\O_\reg(D_X(f))$ which are rational functions on $D_X(f)$).
Explicitly this isomorphism is given by the universal property of localization applied to the restriction $k[X]\to k[D_X(f)]$.

This immediately shows the following two corollaries.

\bcoro

    For $X\subseteq\bA^n$ affine and $f\in k[\bA^n]$, we have
    $$ \O_\reg(D_X(f)) = k[\bA^n]|_{D_X(f)}[1/f] \subseteq \fun(D_X(f),k) , $$
    i.e.\ regular functions on $D_X(f)$ are of the form $p/f$ for $p\in k[\bA^n]$.

\ecoro

\bcoro

    If $(X,\O_X)$ is an affine variety, then for every $f\in\O_X(X)$, the SWF $(D_X(f),\O_X|_{D_X(f)})$ is affine.
    Furthermore, the restriction map $\O_X(X)\to\O_X(D_X(f))$ induces an isomorphism of $k$-algebras
    $$ \O_X(X)_f \xtos\sim \O_X(D_X(f)) . $$

\ecoro

\bnote

    Since $D_X(f)\subseteq X$ is open, the restriction of $\O_X$ to $D_X(f)$ is equal to $\O_X$ but on the subposet of open subsets of $D_X(f)$.

\enote

\bcoro

    An affine variety has a basis consisting of affine varieties.

\ecoro

\bproof

    The $D_X(f)$ form an open basis for $X$ (since they do for $\bA^n$, and $X$ is a subspace and $D_X(f)=D_{\bA^n}(f)\cap X$), and they are affine.
    \qed

\eproof

\bdefn

    Let $X=(X,\O_X)$ be an algebraic variety.
    A {\emphcolor closed/open/locally closed subvariety} of $X$ is an SWF $(Y,\O_X|_Y)$ where $Y\subseteq X$ is closed/open/locally closed.
    A {\emphcolor quasi-affine variety} is an open subvariety of an affine variety.

\edefn

\bcoro

    A closed/open/locally closed subvariety of an algebraic variety is an algebraic variety.

\ecoro

\bproof

    Let $X_i$ be a finite subcovering of $X$ of affine varieties.
    Then $Y=\bigcup_i(Y\cap X_i)$, and to show that $Y$ is an algebraic variety it is sufficient to show that $Y\cap X_i$ is.
    So replacing $X$ by $X_i$ and $Y$ by $Y\cap X_i$, we assume $X$ is affine.

    We consider the cases when $Y$ is open and closed, and this gives us the case when $Y$ is locally closed, since if $Y$ is open in a closed subset
    then it is an open subset of an algebraic variety which is then an algebraic variety.

    If $Y$ is closed, then $(Y,\O_X|_Y)$ is affine since following the isomorphism to a closed subset of $\bA^n$ we see that $Y$ is a closed subset of $\bA^n$
    and thus an affine variety.

    Now if $Y\subseteq X$ is open, then $\set{D_X(f)}_{f\in\O_X(X)}$ is a basis of $X$, and since $X$ is Noetherian, $Y$ is quasi-compact, so it is the finite
    union of $D_X(f)$s.
    Each of these are affine varieties, so $Y$ is an algebraic variety by definition.
    \qed

\eproof

\bexam

    We have shown some quasi-affine varieties are affine, like $D_X(f)$.
    But not all!

    For example let $X=\bA^2-0$, which is quasi affine as it is equal to $D_{\bA^2}(x_1)\cup D_{\bA^2}(x_2)$.
    The inclusion $X\inj\bA^2$ is not an isomorphism, as it is not a homeomorphism, but the induced map $\res\colon k[\bA^2]\to\O_X(X)$ is an isomorphism.
    This cannot happen for affine varieties ($\Gamma$ reflects isomorphisms).

    We see that $X$ is the pushout of $D(x_1),D(x_2)$ along their intersection.
    Thus $\O_X(X)$ is the pullback of $k[x_1,x_2]_{x_1},k[x_1,x_2]_{x_2}$ along $k[x_1,x_2]_{x_1x_2}$.
    That is, $\O_X(X)$ consists of pairs $f_i\in k[x_1,x_2]_{x_i}$ such that $f_1=f_2$ in $k[x_1,x_2]_{x_1x_2}$.
    But a rational function whose denominator is both a power of $x_1$ and a power of $x_2$ must be a polynomial, so $\O_X(X)=k[\bA^2]$.

\eexam


